\section{Benchmark design and experimental protocol}
\label{sec:si-protocol}

\subsection{Task specification and study units}
CloudraBench tests whether agents can construct the requested species and reaction pathways, run the calculations, and draw conclusions supported by the results. It comprises 50 tasks: 20 Easy, 10 Normal, 10 Hard, and 10 Impossible. Easy tasks cover elementary transformations, conformational changes, and basic execution; Normal tasks require identifying intermediates and connecting successive elementary steps; Hard tasks compare competing pathways and stereochemical products; Impossible tasks require recognizing what the allowed electronic-structure methods can support. Each task specifies reactant and target identities, reaction conditions, total charge, spin multiplicity, requested results, and permitted methods. All configurations receive the same specifications for that task. The complete catalogue accompanies the per-task results.

The repeated-task study, P1b, measures repeated runs on four tasks; the full-benchmark study, P2, compares all 50 tasks (Table~\ref{tab:si-study-design}). P1b contains three independent runs for each task and configuration. E1 concerns elimination preorganization; NB1 concerns multistep decarboxylation; HA1 concerns stereochemical ring opening; and IA1 concerns an electronic-structure method limitation. P2 retains one run per task and configuration, so comparisons pair the four configurations on the same task. The 48 repeated-task runs and the 200 full-benchmark runs are analyzed separately.

\begin{table}[htbp]
\centering
\caption{Experimental units in the two studies.}
\label{tab:si-study-design}
\small
\begin{tabular}{p{0.21\textwidth}p{0.31\textwidth}p{0.34\textwidth}}
\toprule
Design & Repeated-task study (P1b) & Full-benchmark study (P2) \\
\midrule
Task set & E1, NB1, HA1, IA1 & All 50 CloudraBench tasks \\
Backbone & DeepSeek-v4.1-flash (1017) & DeepSeek-v4.1-flash (1017); official DeepSeek chat-completions API \\
Configurations & Guided, Standard, Record, Brief-errors & Guided, Standard, Record, Brief-errors \\
Repetitions & Three per task and configuration & One retained run per task and configuration \\
Task-level runs & $4\times4\times3=48$ & $50\times4=200$ \\
Elapsed-time unit & Four-task batch at concurrency two; three batches per configuration & Computational effort on the common 15-task subset \\
\bottomrule
\end{tabular}
\end{table}

\subsection{Workflow configurations and primary comparison}
The primary comparison is Guided against Standard with the same model backbone. Record and Brief-errors test the effects of removing required physical validation and shortening diagnostics, respectively (Table~\ref{tab:si-configurations}). Standard provides computational tools and neutral feedback while leaving physical checks to the agent; explicit delivery-DAG tools are hidden. Guided retains the graph, required scientific checks, and detailed diagnostics. Record retains the tools and records while making physical validation optional. Brief-errors retains guidance and physical validation with shorter diagnostics. Every configuration has the same task and method restrictions.

\begin{table}[htbp]
\centering
\caption{Operational definition of the four workflow configurations. Brief-errors is abbreviated as Brief in figures and tables; Standard is also labelled Backbone.}
\label{tab:si-configurations}
\small
\begin{tabular}{p{0.16\textwidth}p{0.25\textwidth}p{0.25\textwidth}p{0.22\textwidth}}
\toprule
Configuration & Delivery-graph interface & Required physical validation & Diagnostic presentation \\
\midrule
Guided & Explicit & Retained & Detailed \\
Standard & Explicit DAG tools hidden & Removed & Neutral computation feedback \\
Record & Available tools and records retained & Removed & Recording retained \\
Brief-errors & Retained & Retained & Shortened \\
\bottomrule
\end{tabular}
\end{table}

Both studies used DeepSeek-v4.1-flash (1017). All eight saved coordinator and follower configurations and their model-call records use the API alias \texttt{deepseek-v4-flash}.\citep{deepseekv42026} Each run retains its model name, campaign identifier, budget, stopping reason, and calculation files. Cloudra revision 70daa5a identifies the implementation described in the main-text methods.\citep{ref14} Execution, scoring, and offline organization have separate version records; the versions used in each step are taken from its archive.

\subsection{P2 model configuration and campaign versions}
\label{sec:si-p2-model-metadata}

P1b contains four tasks, four workflow configurations, and three runs per combination, giving 48 task runs. P2 contains 50 tasks under the four configurations, giving 200 retained task runs from the P2b and P2c campaigns. Three completed P2b runs were retained (Guided/D3, Record/D1, and Standard/D3); P2c supplied the remaining 197 runs.

The backbone was DeepSeek-v4.1-flash (1017). All eight saved P2 configurations, one per campaign and workflow combination, specify \texttt{deepseek-v4-flash} for both the coordinator and follower. The request-start and completion events in all eight runtime archives record that same identifier. Cloudra used the official DeepSeek endpoint with streaming OpenAI-compatible chat completions, the route selected by \texttt{native\_search=auto} in the campaign implementation. Table~\ref{tab:si-p2-model-metadata} lists the request settings and execution versions, including both the deployment name and the API alias stored in the records.

\begin{table}[htbp]
\centering
\small
\caption{Frozen P2 configuration and campaign metadata. Request settings describe the harness configuration; telemetry counts are provided in the accompanying \texttt{p2\_model\_metadata.json}.}
\label{tab:si-p2-model-metadata}
\begin{tabular}{p{0.27\linewidth}p{0.66\linewidth}}
\toprule
Item & Recorded value \\
\midrule
Deployed model & DeepSeek-v4.1-flash (1017) \\
Provider and API alias & DeepSeek; \texttt{deepseek-v4-flash} in all eight snapshots and their request/completion telemetry \\
Endpoint and transport & \url{https://api.deepseek.com}; streaming chat completions \\
Role slots & Coordinator and follower use the same provider, model identifier, and request settings \\
Reasoning setting & Requested \texttt{reasoning\_effort=minimal} \\
Temperature & 1.0 in the archived run-start settings \\
Context window & 400,000 tokens in the frozen harness configuration \\
Spawned-agent budget & \texttt{max\_tokens=32768} in the archived \texttt{agent.spawned} settings; this budget applies to each agent request. Total run token consumption is a separate quantity. \\
Tool invocation & JavaScript protocol \\
Retained task runs & Three from P2b and 197 from P2c \\
Task concurrency & Two per configuration, eight in total \\
Campaign starts & P2b: 25 October 2026, 18:01:17 UTC+8; P2c: 25 October 2026, 22:18:45 UTC+8 \\
Campaign base & \texttt{53309d0e} \\
Configuration commits & Guided: \texttt{6e89d20c}; Record: \texttt{746c33f0}; Standard: \texttt{c7fc94a1}; Brief-errors: \texttt{3b9831bb}. Both campaigns use these commits. \\
\bottomrule
\end{tabular}
\end{table}

The configuration commits identify the code used to run the experiments. The main-text implementation, later evaluator fixes, and publication-table reproduction have separate version records. The companion metadata file gives the archive path of each configuration and runtime record. Its request-start and completion counts include auxiliary LLM calls and cancelled P2b attempts; the benchmark includes only the 200 retained task runs.

\paragraph{Fixed geometric heuristics.}
Atom maps are ranked with fixed costs: 2 for each formed or broken bond, 1 for each bond-order change, 4 for a hydrogen-count mismatch, 8 for charge or radical-electron mismatches, and 2 for a stereochemical mismatch. Maps within $J_{\min}+2$ are retained when reporting ambiguity. The branch-crowding descriptor uses decay $\rho=0.5$, traversal through depth five, and a congestion flag when the four-branch total $B_{\mathrm{axis}}\geq9.0$. These heuristics rank maps and flag structures for inspection. Physical validation uses the quantum-chemical results. The imaginary-mode participation ratio is compared with 0.5 as an advisory diagnostic; validation requires the first-order saddle and IRC/endpoint calculations.

\paragraph{Conformational sampling.}
Default preorganization places molecular fragments using their topology, relaxes them with guided MMFF94 optimization, and selects candidates by energy and RMSD. For multicomponent systems, it samples 32 relative orientations and retains up to eight guided placements. A separate conformer-search tool uses RDKit by default and offers CREST\citep{crest2020} for additional conformer or noncovalent-complex sampling. The default CREST route uses GFN2-xTB with the quick setting. The resulting geometries are candidates for subsequent calculations and must pass the task's stationary-point and pathway checks.



\subsection{Allowed calculations and reference definitions}
Easy, Normal, and Hard tasks permit xTB, PM7, and B3LYP-D3(BJ)/def2-SVP.\citep{gfn2xtb2019,pm72013,becke1993,lyp1988,stephens1994,d32010,d3bj2011,def22005} Impossible tasks are restricted to xTB and PM7. Computational quantities and termination status are read from native ORCA, Gaussian, and MOPAC outputs.\citep{orca2025,gaussian16,mopac2026,ref11} Each result retains its program, calculation level, species, and task conditions.

The reference answer (Gold) specifies the mechanistic species and elementary steps, stereochemical identities, applicable numerical targets, and explicit negative targets. Numerical activation barriers use electronic energies at a consistent level of theory:
\begin{equation}
\Delta E_e^{\ddagger}=E_{\mathrm{TS},e}-E_{\mathrm{ref},e}.
\label{eq:si-barrier}
\end{equation}
These benchmark barriers use electronic energies alone. The energy reference is the verified stable minimum of the corresponding species. Declared conformational connections within the same species can share that stable reference; a conformational conversion itself uses its own starting conformer. Distinct stereochemical substrates and independent channel contact complexes retain separate references. The IRC endpoint and the energy reference are recorded separately, since they can be different conformers of the same species.

References follow the individual task definitions. NB1\_E2 and NQ5\_E2 are exempt from numerical references under the specified proton-transfer policy. IB9 combines a numerical first-step reference of 5.65~kcal~mol$^{-1}$ with a negative second-step target. Section~\ref{sec:si-scoring} specifies the qualitative selectivity targets and numerical tolerances. Each reference record separately states whether the task and Gold graph are available and whether the complete calculation chain has been physically validated.

\subsection{Physical evidence associated with a mechanism}
A mechanism graph represents species occurrences as nodes and elementary steps as edges. A species minimum requires a converged optimization and zero imaginary frequencies. An accessible transition-state step requires a first-order saddle with exactly one imaginary frequency associated with the intended reaction, plus IRC results connecting that saddle to the declared endpoint species.\citep{ref09,ref10} Normal termination, geometry convergence, frequencies, and pathway connectivity are checked separately.

Both IRC directions must start at the same transition state. The actual outer frame from each direction is then optimized to a minimum with zero imaginary frequencies and the declared chemical identity and stereochemistry. This outer-frame--optimization--frequency chain can establish the connection from a short IRC. The frames and links between calculations are archived. Native calculation results establish the endpoint connection; geometric interval measurements aid interpretation. For separated products, the required constituent minima are verified; an intact molecular complex is evaluated as a whole-system species.

These requirements apply to the species and steps requested by the task. A negative conclusion marks the target as physically inaccessible under the permitted methods and records the reason. A species-only negative question is evaluated directly on that species claim.

\subsection{Offline organization and publication records}
Formal evaluation represents the research completed by all four configurations in the same mechanism-graph format. Human organizers read each run's conversations, tool records, reports, geometries, and native calculation outputs, then encode the agent's stated species, steps, mechanistic intent, and conclusions. They recover graph entries, repair file references, and convert quantities already present in the run; each recovered field records its source. The evaluator searches that run's sandbox for native calculation results and applies the same physical criteria and scoring rules to every configuration. Organization is confined to that run's existing calculations and stated conclusions; Gold supplies the evaluation reference. The procedure and field sources allow the graph to be reconstructed from the original run materials.

The full-benchmark table contains 200 runs, all 600 Graph, Self-consistency, and Accuracy values, and their evaluation sources. The weighted components reproduce the final totals to four decimal places; Section~S7 specifies the references and evaluator revision. The tables also record whether the run wrote a root mechanism graph and whether it was organized offline. Readers can therefore inspect graph delivery separately from the score for the run's research results.

The repeated-task study retains its historical results for all four configurations. The later offline organization updates Standard separately; the other configurations in the 48-run study retain their historical scores. Analyses of this archive identify which version supplies each result.

\subsection{Paired statistical analysis and time measurements}
For a task $i$ and configuration $a$, let $S_{ia}$ denote its composite score. Configuration means give each task equal weight within the reported set. The primary paired difference is
\begin{equation}
\overline{\Delta}_{\mathrm{G-S}}=\frac{1}{n}\sum_{i=1}^{n}
\left(S_{i,\mathrm{Guided}}-S_{i,\mathrm{Standard}}\right).
\label{eq:si-paired-difference}
\end{equation}
The same calculation is used within each tier and for the combined Normal/Hard subset. Weighted component contributions are the mean paired component differences multiplied by 0.1, 0.3, and 0.6, respectively. They add to the composite difference up to the reporting precision of the final scores.

The 95\% percentile intervals use 100,000 task-cluster bootstrap samples with seed 20261031. Each draw selects a task with its four configuration records together, preserving the paired comparison.

Elapsed time in the repeated-task study runs from the recorded start to finish of each four-task batch at concurrency two: three batches per configuration, or 12 batches in total. It is measured in batch wall-clock hours. Full-benchmark effort is measured on the 15 tasks with records for all four configurations, comprising 60 task-level runs. Calculation counts and cumulative quantum-chemical job-hours are reported; job-hours sum the individual calculation durations within each run. For Standard/M1, native-log CPU time is divided by four to estimate job duration. The tables keep job-hours separate from core-hours and from the elapsed time of an entire P2 research run.
